\section{Grupa wolna: definicja, konstrukcja i przykłady}
\subsection{Podstawowe pojęcia z teorii grup}
\begin{definicja}[Grupa]
Grupa to zbiór $G$ z działaniem $(g,h)\mapsto gh$, które jest łączne, ma element neutralny $e$ i dla każdego $g\in G$ ma element odwrotny $g^{-1}$. Oznacza to $(gh)k=g(hk)$, $eg=ge=g$, $gg^{-1}=g^{-1}g=e$. Grupa jest \emph{przemienna}, gdy dodatkowo $gh=hg$ dla wszystkich $g,h$. W tym skrypcie przemienności na ogół nie zakładamy.
\end{definicja}
\begin{przyklad}
$\Z$ z dodawaniem jest grupą: elementem neutralnym jest $0$, a odwrotnością $n$ jest $-n$. Bijekcje danego zbioru tworzą grupę ze składaniem funkcji. W tym drugim przykładzie kolejność może mieć znaczenie. Dla funkcji przyjmujemy $(fg)(x)=f(g(x))$: najpierw działa czynnik z prawej strony.
\end{przyklad}
\begin{definicja}[Homomorfizm, podgrupa, generatory]
Homomorfizm $\varphi:G\to H$ spełnia $\varphi(gh)=\varphi(g)\varphi(h)$. Jeżeli jest różnowartościowy, nazywamy go \emph{monomorfizmem}. Podgrupa $K\le G$ to podzbiór zawierający $e$ i zamknięty na iloczyny i odwrotności. Mówimy, że $a,b$ generują $G$, gdy każdy element $G$ jest skończonym iloczynem elementów $a,a^{-1},b,b^{-1}$; dopuszczamy iloczyn pusty, równy $e$. Piszemy wtedy $\langle a, b\rangle=G.$
\end{definicja}
Homomorfizm zachowuje element neutralny i odwrotności. Na przykład $\varphi(e)=\varphi(e)^2$; pomnożenie przez $\varphi(e)^{-1}$ daje $\varphi(e)=e_H$. Następnie $\varphi(g)\varphi(g^{-1})=\varphi(e)=e_H$.

\subsection{Definicja grupy wolnej}
\begin{definicja}[Grupa wolna na dwóch generatorach]
\label{def:grupa-wolna}
Elementy $a,b$ grupy $F$ są \emph{wolnymi generatorami},
jeżeli generują $F$ i spełniają następującą własność: dla każdej grupy
$H$ i dowolnych elementów $u,v\in H$ istnieje dokładnie jeden homomorfizm
$\varphi:F\to H$ taki, że $\varphi(a)=u$ oraz $\varphi(b)=v$.
Taką grupę oznaczamy $F(a,b)$ lub $F_2$. Własność tę nazywamy
\emph{własnością uniwersalną}.
\end{definicja}

Znaczenie definicji jest następujące: obrazy $a$ i $b$ można wybrać
zupełnie dowolnie, bez konieczności sprawdzania dodatkowych relacji.
Gdybyśmy narzucili $a^2=e$, to obraz $u$ musiałby spełniać $u^2=e_H$;
gdybyśmy narzucili $ab=ba$, to obrazy musiałyby komutować.
Grupa wolna nie narzuca żadnych relacji poza wynikającymi
z aksjomatów grupy.

Nie wymagamy przy tym, aby $u,v$ generowały całą grupę $H$.
Ponieważ $a,b$ generują $F$, obraz homomorfizmu jest dokładnie
podgrupą $\langle u,v\rangle\leq H$. Homomorfizm $\varphi$
jest więc surjektywny wtedy i tylko wtedy, gdy
$H=\langle u,v\rangle$.

Równoważnie można ograniczyć własność uniwersalną do grup
generowanych przez dwa wyróżnione elementy. Istotnie, dla dowolnych
$u,v\in H$ stosujemy ją do $K=\langle u,v\rangle$, otrzymując
homomorfizm $\widetilde\varphi:F\to K$, a następnie składamy go
z inkluzją $K\hookrightarrow H$. Jednoznaczność wynika z faktu,
że homomorfizm z $F$ jest całkowicie wyznaczony przez obrazy
generatorów $a,b$.

Ta intuicja nie zastępuje konstrukcji. Teraz zbudujemy grupę, sprawdzimy wszystkie aksjomaty i udowodnimy powyższą własność.

\subsection{Litery, słowa i redukcja}
Bierzemy cztery formalne symbole $\mathcal A=\{a,\bar a,b,\bar b\}$. Na razie kreska nie oznacza obliczonej odwrotności: ustalamy tylko parowanie liter $a\leftrightarrow\bar a$ oraz $b\leftrightarrow\bar b$. Dla litery $s$ oznaczamy przypisaną jej literę odwrotną przez $s^{-1}$. Później $\bar a$ rzeczywiście będzie odwrotnością elementu $a$.
\begin{definicja}[Słowo zredukowane]
Słowem jest skończony ciąg liter $s_1\cdots s_n$ z $\mathcal A$. Słowo puste oznaczamy $e$; jego długość wynosi zero. Słowo jest \emph{zredukowane}, jeśli nie zawiera sąsiedniej pary liter $ss^{-1}$. Redukcja elementarna usuwa taką parę.
\end{definicja}
\begin{przyklad}
Słowo $a\bar a b\bar b ba$ redukuje się następująco: $a\bar a b\bar b ba\longrightarrow b\bar b ba\longrightarrow ba$. Słowo $ab\bar a\bar b$ jest już zredukowane. Nie wolno skrócić liter $a$ i $\bar a$, ponieważ nie sąsiadują. W grupie wolnej nie wolno przestawić liter, aby ułatwić sobie skracanie.
\end{przyklad}

\subsection{Dlaczego wynik redukcji jest jednoznaczny?}
To istotny punkt: bez jednoznaczności nie mielibyśmy dobrze określonego działania. Nie wystarcza powiedzieć „skracamy, ile się da”.
\begin{definicja}[Algorytm stosu]
Czytamy słowo od lewej do prawej, przechowując pomocnicze słowo $r$, początkowo puste. Po przeczytaniu litery $s$:
\begin{itemize}
\item jeśli ostatnią literą $r$ jest $s^{-1}$, usuwamy ją;
\item w przeciwnym razie dopisujemy $s$ na końcu $r$.
\end{itemize}
Wynik po przeczytaniu całego słowa $w$ oznaczamy $\red(w)$. Pomocnicze słowo jest zredukowane po każdym kroku.
\end{definicja}
\begin{lemat}[Odczytanie pary odwrotnej niczego nie zmienia]\label{lem:stos}
Jeśli przed odczytaniem liter $s,s^{-1}$ na stosie jest słowo zredukowane $r$, to po odczytaniu obu liter na stosie znów jest $r$.
\end{lemat}
\begin{proof}
Rozważamy dwa przypadki. Jeśli $r$ nie kończy się literą $s^{-1}$, pierwsza operacja dopisuje $s$, a druga natychmiast usuwa dopisane $s$. Wynik to $r$. Jeśli $r=us^{-1}$, pierwsza operacja pozostawia $u$. Ponieważ $r$ było zredukowane, $u$ nie kończy się literą $s$. Druga operacja dopisuje zatem $s^{-1}$, a nie kasuje końca $u$. Ponownie otrzymujemy $us^{-1}=r$. Działa to również dla pustego $u$.
\end{proof}
\begin{twierdzenie}[Jednoznaczność redukcji]\label{tw:redukcja}
Każde słowo $w$ można przez kolejne redukcje elementarne sprowadzić do dokładnie jednego słowa zredukowanego, którym jest $\red(w)$.
\end{twierdzenie}
\begin{proof}
\emph{Istnienie.} Każde usunięcie pary zmniejsza długość o $2$, więc ciąg usunięć musi się skończyć. Wynik końcowy jest zredukowany. Algorytm stosu również można zrealizować przez takie usunięcia: redukujemy przeczytany prefiks, dopisujemy następną literę i ewentualnie usuwamy jedną końcową parę. Indukcja po długości prefiksu uzasadnia, że końcowy stos jest osiągalnym zredukowanym wynikiem.

\emph{Niezależność od wyboru par.} Słowa $pss^{-1}q$ i $pq$ mają ten sam wynik algorytmu stosu. Po prefiksie $p$ stos jest ten sam; dwie dodatkowe litery nie zmieniają go na mocy lematu~\ref{lem:stos}; następnie czytamy ten sam sufiks $q$.

Zatem każde usunięcie pary zachowuje $\red(w)$. Jeżeli pewien ciąg usunięć kończy się słowem zredukowanym $z$, to $\red(w)=\red(z)=z$. Ostatnia równość zachodzi, ponieważ algorytm stosu nie skraca słowa, w którym nie ma odwrotnych sąsiadów. Każdy wynik końcowy jest więc taki sam.
\end{proof}
\begin{wniosek}\label{wn:redukcja-iloczyn}
Dla dowolnych słów $u,v$ mamy $\red(\red(u)v)=\red(uv)$ oraz $\red(u\red(v))=\red(uv)$.
\end{wniosek}
\begin{proof}
Redukcje prowadzące od $u$ do $\red(u)$ można wykonać wewnątrz początkowego fragmentu $uv$. Analogicznie redukcje $v$ można wykonać wewnątrz końcowego fragmentu. Twierdzenie~\ref{tw:redukcja} mówi, że ostateczny wynik od tego nie zależy.
\end{proof}

\subsection{Mnożenie słów zredukowanych}
Niech $F_2$ będzie zbiorem wszystkich słów zredukowanych. Definiujemy $u\cdot v=\red(uv)$, gdzie $uv$ po prawej oznacza złączenie słów.
\begin{stwierdzenie}\label{st:grupa-slow}
Tak zdefiniowane mnożenie czyni $F_2$ grupą. Elementem neutralnym jest słowo puste $e$, a odwrotnością $s_1\cdots s_n$ jest $s_n^{-1}\cdots s_1^{-1}$.
\end{stwierdzenie}
\begin{proof}
Zamkniętość wynika z tego, że $\red$ daje słowo zredukowane. Łączność nie jest dodatkowym założeniem: z wniosku~\ref{wn:redukcja-iloczyn} otrzymujemy $(u\cdot v)\cdot w=\red(\red(uv)w)=\red(uvw)$ oraz $u\cdot(v\cdot w)=\red(u\red(vw))=\red(uvw)$. Doklejenie słowa pustego nic nie zmienia, zatem jest ono neutralne.

Słowo $s_n^{-1}\cdots s_1^{-1}$ jest zredukowane: sąsiednia para odwrotna w nim odpowiadałaby sąsiedniej parze odwrotnej w pierwotnym słowie. W iloczynie $(s_1\cdots s_n)(s_n^{-1}\cdots s_1^{-1})$ kasujemy kolejno pary od środka na zewnątrz i otrzymujemy $e$. Tak samo dzieje się w odwrotnej kolejności iloczynu.
\end{proof}
\begin{twierdzenie}[Własność uniwersalna konstrukcji]\label{tw:uniwersalnosc}
Zbudowana grupa $F_2$ jest grupą wolną na $a,b$ w sensie definicji~\ref{def:grupa-wolna}.
\end{twierdzenie}
\begin{proof}
Weźmy dowolną grupę $H$ i $u,v\in H$. Zastąpmy litery $a,\bar a,b,\bar b$ odpowiednio elementami $u,u^{-1},v,v^{-1}$, a słowo puste elementem $e_H$. Dla dowolnego słowa $w$ niech $V(w)$ oznacza iloczyn przypisanych mu elementów $H$ w tej samej kolejności. Każda usuwana para $ss^{-1}$ ma iloczyn $e_H$, więc $V(\red(w))=V(w)$. Ograniczenie $V$ do słów zredukowanych oznaczmy przez $\varphi$. Dla zredukowanych słów $w,z$ mamy $\varphi(\red(wz))=V(\red(wz))=V(wz)=V(w)V(z)=\varphi(w)\varphi(z)$. To homomorfizm o żądanych obrazach $a,b$. Jednoznaczność wynika z tego, że homomorfizm musi zachowywać iloczyny i odwrotności: jego wartość na każdym słowie jest wymuszona przez obrazy liter.
\end{proof}
\begin{uwaga}[Jednoznaczność grupy wolnej]
Jeśli dwie grupy spełniają własność uniwersalną dla dwóch wskazanych generatorów, otrzymujemy homomorfizmy w obie strony zachowujące generatory. Ich złożenia też zachowują generatory, tak samo jak odwzorowanie identycznościowe. Z jednoznaczności we własności uniwersalnej złożenia są identycznościami. Grupy są więc izomorficzne.
\end{uwaga}

\subsection{Przykłady i kontrprzykłady}
\begin{przyklad}[Jeden generator]
Przy alfabecie $\{a,\bar a\}$ każde niepuste słowo zredukowane składa się wyłącznie z $a$ albo wyłącznie z $\bar a$: zmiana rodzaju litery tworzyłaby parę odwrotną. Dostajemy elementy $a^n$, $n\in\Z$, i regułę $a^ma^n=a^{m+n}$. Grupa wolna na jednym generatorze jest więc izomorficzna z $(\Z,+)$. Jest przemienna. Wolność nie oznacza automatycznie nieprzemienności.
\end{przyklad}
\begin{przyklad}[Dwa generatory]
W $F_2$ słowa $ab$ i $ba$ są różnymi słowami zredukowanymi, więc $ab\ne ba$. W szczególności $aba^{-1}b^{-1}\ne e$. Jest to przykład komutatora. Niezerowa długość zredukowanego słowa jest tutaj pełnym dowodem, że nie otrzymaliśmy elementu neutralnego. Na przykład $(ab)^{-1}=b^{-1}a^{-1}$, a nie $a^{-1}b^{-1}$. Iloczyn $ab\cdot b^{-1}a^{-1}$ redukuje się do $e$.
\end{przyklad}
\begin{przyklad}[Grupa $\Z^2$ nie jest wolna na swoich dwóch standardowych generatorach]
Dla $u=(1,0)$ i $v=(0,1)$ zachodzi $u+v=v+u$. Zastąpienie $a$ przez $u$, a $b$ przez $v$ daje homomorfizm $F_2\to\Z^2$, ale komutator przechodzi na $(0,0)$. Homomorfizm nie jest różnowartościowy. W $\Z^2$ pamiętamy sumy wykładników, lecz tracimy kolejność liter; w $F_2$ ta kolejność ma znaczenie.
\end{przyklad}
\begin{przyklad}[Nieprzemienność nie wystarcza]
Permutacje $(12)$ i $(23)$ zbioru $\{1,2,3\}$ nie komutują, ale $(12)^2=e$. Nie mogą być wolnymi generatorami. Podobnie wybranie dwóch niekomutujących obrotów w przestrzeni nie dowodzi jeszcze, że generują $F_2$: trzeba wykluczyć wszystkie niepuste słowa zredukowane, nie tylko jeden komutator.
\end{przyklad}
\begin{stwierdzenie}[Przeliczalność]\label{st:przeliczalnosc-f2}
Grupa $F_2$ jest przeliczalnie nieskończona. Słów zredukowanych o długości $n\ge1$ jest dokładnie $4\cdot3^{n-1}$.
\end{stwierdzenie}
\begin{proof}
Pierwszą literę wybieramy na cztery sposoby. Każda następna może być dowolna poza odwrotnością poprzedniej, więc mamy trzy możliwości. Możemy wypisać słowa według długości, a w ramach jednej długości według ustalonego porządku liter. To daje wyliczenie wszystkich elementów. Jest ich nieskończenie wiele, bo $e,a,a^2,\ldots$ są różne.
\end{proof}

